\section*{Overview}

Suffix array is the static full-string structure I want when the problem is about lexicographic order of suffixes,
substring comparison, or repeated substrings, but I still want something lighter than a suffix tree.

\section*{Problem-Driven Motivation}

Suppose a problem asks for one of the following on a string of length \(n \le 2 \cdot 10^5\):

\begin{itemize}
  \item count distinct substrings,
  \item find the lexicographically smallest occurrence of a pattern,
  \item analyze repeated substrings,
  \item compare many substrings lexicographically.
\end{itemize}

The naive view is to work directly with substrings, but there are \(O(n^2)\) of them. The structural observation is
that every substring is a prefix of some suffix. Once all suffixes are sorted, a lot of substring logic becomes local.

\section*{Recognition Pattern}

Suffix array is a strong fit when:

\begin{itemize}
  \item the string is static,
  \item lexicographic suffix order matters,
  \item pattern search by binary search is acceptable,
  \item repeated-substring questions reduce to common prefixes of nearby suffixes.
\end{itemize}

Signals that usually point here:

\begin{itemize}
  \item ``sort suffixes / rotations'',
  \item ``count distinct substrings'',
  \item ``longest repeated substring'',
  \item ``pattern queries on one fixed text''.
\end{itemize}

\section*{Derivation}

The suffix array \(\texttt{sa}\) stores all suffix starting positions in lexicographic order.

The hard part is building it without comparing long suffix strings directly. The doubling algorithm solves this by
ranking prefixes of increasing length.

\paragraph{Round \(k\).}
Assume we already know the lexicographic classes of all length-\(2^k\) prefixes. Then the first \(2^{k+1}\) characters
of suffix \(i\) are summarized by the pair
\[
(\texttt{cls}[i], \texttt{cls}[i + 2^k]).
\]

So sorting suffixes by their first \(2^{k+1}\) characters becomes sorting pairs of small integers. After enough rounds,
the compared length exceeds the string length, so the full suffix order is known.

\paragraph{Why LCP appears naturally.}
Once suffixes are sorted, repeated-substring questions are local because two suffixes with a long common prefix must sit
next to each other in suffix-array order. That is why the LCP array is almost always built together with \(\texttt{sa}\).

\section*{Worked Problem}

\paragraph{Problem.}
Count the number of distinct substrings of a string \(s\).

\paragraph{Why naive fails.}
There are \(O(n^2)\) substrings, and even hashing all of them is usually too expensive at \(n = 2 \cdot 10^5\).

\paragraph{Key observation.}
Suffix \(s[sa[i] \dots n-1]\) contributes all of its prefixes as substrings. That is
\[
n - sa[i]
\]
substrings in total. But the first \(\texttt{lcp}[i-1]\) of them were already seen in the previous suffix.

\paragraph{Therefore}
\[
\text{distinct substrings}
=
\sum_i (n - sa[i]) - \sum_i lcp[i].
\]

\paragraph{Why only adjacent LCPs matter.}
In sorted order, the largest already-seen common prefix of the current suffix is always realized by one of its
neighbors, so subtracting adjacent LCPs removes exactly the duplicates introduced by ordering.

\section*{Implementation Reasoning}

The implementation has three moving parts:

\begin{itemize}
  \item append one unique terminator smaller than all real characters,
  \item build \(\texttt{sa}\) with the doubling method,
  \item build \(\texttt{lcp}\) with Kasai's linear scan.
\end{itemize}

Two practical details matter a lot:

\begin{itemize}
  \item the terminator ensures no suffix is a prefix of another without a clean ordering,
  \item \(\texttt{lcp[i]}\) belongs to adjacent suffixes \(\texttt{sa[i]}\) and \(\texttt{sa[i+1]}\), not to original
  string positions.
\end{itemize}

The code returns both arrays and includes a distinct-substring application function.

\section*{Correctness Intuition}

The doubling invariant is the real proof:

\begin{quote}
After round \(k\), the class array correctly ranks all prefixes of length \(2^k\).
\end{quote}

Sorting pairs of those classes therefore sorts prefixes of length \(2^{k+1}\). Once the compared length covers the
whole suffixes, the order is final.

Kasai's algorithm is linear because its matched prefix length only decreases by at most one between consecutive starting
positions.

\section*{Complexity Analysis}

\begin{itemize}
  \item suffix array by doubling: \(O(n \log n)\),
  \item LCP by Kasai: \(O(n)\),
  \item memory: \(O(n)\).
\end{itemize}

\section*{Common Pitfalls}

\begin{itemize}
  \item Forgetting the unique terminator.
  \item Misaligning \(\texttt{lcp}\) with suffix-array order.
  \item Using suffix array when the real need is online substring extension, which is more natural for suffix automaton.
  \item Assuming arbitrary substring LCP queries are free without adding an RMQ layer over \(\texttt{lcp}\).
\end{itemize}

\section*{Variants and Failure Modes}

\begin{itemize}
  \item Add RMQ over \(\texttt{lcp}\) for fast LCP queries between arbitrary suffixes.
  \item Use suffix automaton when the string grows online or substring occurrence aggregation is more natural.
  \item SA-IS gives linear construction, but the contest constant and implementation complexity are much higher.
\end{itemize}

\section*{Practice Problems}

\begin{itemize}
  \item Count distinct substrings.
  \item Find the longest repeated substring.
  \item Pattern existence and lexicographic substring queries on a fixed string.
\end{itemize}

\section*{References}

\begin{itemize}
  \item \href{https://cp-algorithms.com/string/suffix-array.html}{cp-algorithms: Suffix Array}.
  \item \href{https://usaco.guide/adv/suffix-array?lang=cpp}{USACO Guide: Suffix Array}.
  \item \href{https://codeforces.com/catalog}{Codeforces Catalog}.
\end{itemize}
